\section{Introduction}

Closed-circuit television (CCTV) is one of the most widely deployed sensing infrastructures in modern society. Campuses, roadways, industrial facilities, retail premises, and public spaces depend on continuous video capture for deterrence, situational awareness, incident response, and post-event investigation \cite{Shao2018Smart,Chen2019DIVS}. As camera counts, resolutions, and retention requirements grow, the cost of storing and transmitting surveillance video grows with them. A single 1080p camera operating around the clock produces tens of gigabytes per day; a multi-camera installation quickly reaches storage and bandwidth demands that dominate the total cost of ownership of the surveillance system. Because most surveillance footage contains long uneventful stretches of static background, there is a strong practical incentive to encode video content-adaptively: to spend bits where they matter for security purposes and to economize where they do not \cite{Xiao2017Sensitive,Wu2020FGBG}.

Conventional video coding standards such as H.264/AVC \cite{Wiegand2003H264} and H.265/HEVC \cite{Sullivan2012HEVC} are highly efficient general-purpose compressors, but in their default uniform configuration they distribute encoding quality evenly across the frame. Foreground objects of security interest---people, vehicles, carried objects---receive the same treatment as empty pavement, walls, or sky. Region of Interest (ROI) coding addresses this mismatch by assigning higher quality to selected spatial regions and lower quality elsewhere \cite{Huang2009MultiPriority,Xue2016FastROI}. When the ROI is chosen dynamically from live scene understanding rather than fixed in advance, the encoder can track activity as it moves through the scene. This thesis is concerned with exactly such a system: an adaptive CCTV compression pipeline in which real-time object detection drives dynamic ROI control, and ROI control in turn drives compression, buffering, and visualization decisions.

\section{Problem Statement}

The central problem addressed by this work is the inefficient use of bitrate and storage in continuous CCTV systems that encode all scene regions uniformly, regardless of their security relevance. Three consequences follow. First, at constrained bitrates, uniform encoding degrades the very regions---faces, persons, vehicles, actions---that operators and investigators need most. Second, at the quality levels required to keep those regions usable, the system pays for full-frame fidelity that is unnecessary for static background. Third, naive quality reduction (for example, aggressive global quantization) damages forensic usability without any principled notion of which content must be preserved. Existing ROI-based methods demonstrate that content-aware allocation can reduce bandwidth substantially \cite{Wu2020FGBG,Wang2023VaBUS}, but reported gains are obtained under specific datasets, codecs, and ROI constructions, and they do not automatically transfer to a new deployment with different cameras, scenes, and computational constraints. There is therefore a need for a deployable, edge-compatible prototype whose behaviour is measured honestly under controlled, reproducible conditions---including the possibility that a given ROI construction fails to beat uniform encoding.

\section{Motivation}

The motivation for this work is both economic and operational. Economically, bandwidth and storage dominate CCTV operating cost, particularly for multi-camera installations and for deployments in bandwidth-constrained environments. Technically, the convergence of three mature ingredients makes dynamic ROI control practical at the edge: (i) efficient real-time object detectors in the YOLO family; (ii) flexible software encoders such as those provided by FFmpeg, including HEVC; and (iii) inexpensive edge-cloud software stacks (Python services with web dashboards) capable of carrying detection metadata alongside video. Prior work has shown that edge-based processing can substantially reduce transmission requirements \cite{Nguyen2020Scalable,Li2022EdgeComp} and that semantic or recognition-aware allocation can preserve analytic performance at lower bitrates \cite{Beye2023Recognition,Elgamal2020SIEVE}. The research direction toward video coding for machines \cite{Duan2020VCM} further supports the view that surveillance compression should be judged by the preservation of task-relevant content rather than by global fidelity alone. This thesis builds a concrete system in that spirit and subjects it to disciplined evaluation.

\section{Objectives}

\subsection{Main Objective}

The main objective of this work is to design, implement, and evaluate an adaptive CCTV compression system that uses dynamic ROI control---driven by real-time object detection and a heuristic scene-risk assessment---to allocate encoding resources across an edge-cloud pipeline.

\subsection{Specific Objectives}

\begin{itemize}
\item To implement an edge node that captures video, detects objects with YOLOv8n, generates prioritized ROI maps with temporal persistence, assesses scene risk (NORMAL/ALERT/CRITICAL), and encodes adaptively with FFmpeg.
\item To implement a Lagrangian-inspired bitrate allocation strategy that distributes a bitrate budget between ROI and background regions according to ROI importance and system state.
\item To implement tamper-evident event buffering with pre-event and post-event context and SHA-256 hash-chained integrity manifests.
\item To implement a central server and a Next.js dashboard for live monitoring, metrics, heatmaps, playback, and privacy modes.
\item To probe the behaviour of the detection and risk-scoring components on public datasets (CDnet-2014 and UCF-Crime).
\item To calibrate the heuristic risk score against manually verified per-second labels using a group-aware (Leave-One-Video-Out) protocol and to report its limits.
\item To verify the tamper-evidence properties of the forensic chain under controlled original, clean-control, recompressed, and tampered conditions.
\item To compare the proposed adaptive ROI compression against uniform encoding and alternative ROI baselines under a locked same-codec protocol, and to report the outcome---favourable or not---without distortion.
\end{itemize}

\section{Innovation of Project/Research}

The innovative aspect of this work lies in the integration, under one deployable prototype, of four elements that are more commonly studied in isolation: (i) YOLOv8n-driven dynamic ROI generation with temporal smoothing (ROI TTL and box merging); (ii) a heuristic risk score with hysteresis-guarded NORMAL/ALERT/CRITICAL states that jointly governs compression strength, resolution, GOP structure, and event buffering; (iii) a Lagrangian-inspired ROI/background bitrate allocator implemented as a practical heuristic rather than a claimed global optimizer; and (iv) a tamper-evident forensic chain bound to the event-buffering path. The evaluation methodology is itself part of the contribution: public-dataset probing, group-aware calibration with explicit degeneracy analysis, controlled forensic tests with pre-registered expectations, and a locked same-codec baseline comparison with shared detection caches and withheld BD-rate where interpolation preconditions fail.

\section{Significance of Project/Research}

For practitioners, the thesis provides a reproducible edge-cloud CCTV prototype with documented parameters, configuration profiles, and failure modes. For researchers, it provides a cautionary, fully documented negative result: the tested ROI construction did not outperform uniform encoding at matched bitrate on the evaluated sequence, which clarifies the conditions under which ROI-based gains reported in the literature \cite{Wu2024ROIDVC,Wu2025How2Compress,Wang2025CAVC} can and cannot be expected. For operators and policymakers, the forensic-integrity experiments delineate what a hash-chained event buffer can demonstrate (exact-byte tamper evidence) from what it cannot (legal admissibility or certification), consistent with the chain-of-evidence concerns raised in prior work \cite{Kanwal2020Chain}. The thesis therefore contributes usable engineering together with the evidence needed to judge that engineering honestly.

\section{Organization of Book Chapters}

The remainder of this book is organized as follows. Chapter~2 reviews the literature across 45 references, covering video coding standards, ROI-based compression, edge-cloud analytics, adaptive allocation, forensic integrity, and anomaly detection, and identifies the research gap. Chapter~3 documents the modern hardware and software tools used. Chapter~4 presents the methodology and design approach, including the conceptual pipeline, ROI policy, risk scoring, Lagrangian-inspired allocation, and evaluation protocols. Chapter~5 details the system architecture, and Chapter~6 the implementation. Chapter~7 reports and discusses the results of all four evaluation phases. Chapter~8 describes the web application (central server and dashboard). Chapter~9 covers project and time management. Chapters~10--12 address ethical considerations; system performance, sustainability and societal impact; and societal, health, safety, legal, and cultural aspects. Chapter~13 reflects on lifelong learning and self-development, and Chapter~14 concludes with contributions, limitations, and future work.
